\documentclass[10pt,twocolumn,letterpaper]{article}
\usepackage[rebuttal]{cvpr}

\usepackage{graphicx}
\usepackage{amsmath}
\usepackage{amssymb}
\usepackage{booktabs}
\usepackage[table]{xcolor}
\usepackage{multirow}
\usepackage{array}
\usepackage{pgf}
\usepackage{microtype}
\usepackage[labelsep=period]{caption}
\captionsetup{font=small}
\captionsetup[table]{aboveskip=0pt, belowskip=0pt}
\captionsetup[figure]{aboveskip=2pt, belowskip=0pt}

\usepackage[pagebackref,breaklinks,colorlinks,bookmarks=false]{hyperref}
\usepackage[capitalize]{cleveref}
\crefname{section}{Sec.}{Secs.}
\Crefname{section}{Section}{Sections}
\Crefname{table}{Table}{Tables}
\crefname{table}{Tab.}{Tabs.}

% 表格命令
\newcommand{\mapcellS}[2]{\fontsize{2.2pt}{3pt}\selectfont#1\if\relax\detokenize{#2}\relax\else\makebox[1.2em][l]{\hspace{0.12em}\raisebox{-0.15ex}{\fontsize{1.5pt}{1.8pt}\selectfont\textcolor{black}{(}\fontsize{1.5pt}{1.8pt}\selectfont\textcolor{red}{+#2}\textcolor{black}{)}}}\fi}
\newcommand{\effcell}[3]{\fontsize{2.2pt}{3pt}\selectfont#2\raisebox{-0.15ex}{\fontsize{1.5pt}{1.8pt}\selectfont\textcolor{black}{(}\scalebox{0.55}{\textcolor{green}{$\downarrow$}}\fontsize{1.5pt}{1.8pt}\selectfont\textcolor{green}{#3\%}\fontsize{1.5pt}{1.8pt}\selectfont\textcolor{black}{)}}}

\newcommand{\fpscell}[1]{\fontsize{2.2pt}{3pt}\selectfont#1}

\definecolor{impBest}{HTML}{FFE0C2}
\definecolor{impSecond}{HTML}{FFEDDA}
\definecolor{impThird}{HTML}{FFF6EC}
\definecolor{impBase}{HTML}{FFA500}
\colorlet{impStep1}{impBase!24!white}
\colorlet{impStep2}{impBase!40!white}
\colorlet{impStep3}{impBase!56!white}
\colorlet{impStep4}{impBase!72!white}
\colorlet{impStep5}{impBase!86!white}
\colorlet{impStep6}{impBase!96!white}

\newcommand{\applyimpcolor}[1]{%
\ifcase#1\relax
\or\cellcolor{impStep1}\or\cellcolor{impStep2}\or\cellcolor{impStep3}\or\cellcolor{impStep4}\or\cellcolor{impStep5}\or\cellcolor{impStep6}%
\else\cellcolor{impStep6}%
\fi
}

\newcommand{\mapcellSI}[2]{%
\if\relax\detokenize{#2}\relax
  \mapcellS{#1}{#2}%
\else
  \ifdim#2pt<0.005pt
    \textbf{\mapcellS{#1}{#2}}%
  \else
    \pgfmathtruncatemacro{\impsteps}{min(6,max(1,floor(#2/0.005)))}%
    \applyimpcolor{\impsteps}\textbf{\mapcellS{#1}{#2}}%
  \fi
\fi
}

\newcommand{\tidedown}{\raisebox{-0.15ex}{\scalebox{0.8}{\textcolor{black}{(}\textcolor{green}{$\downarrow$}\textcolor{black}{)}}}}
\newcommand{\tideupAB}[2]{\raisebox{-0.15ex}{\fontsize{3pt}{3.5pt}\selectfont\textcolor{black}{(}\textcolor{red}{+#1}\textcolor{black}{)}}}

% 定义Reviewer颜色
\definecolor{rev6color}{RGB}{153, 102, 0}
\definecolor{rev7color}{RGB}{0, 153, 153}
\definecolor{rev8color}{RGB}{0, 153, 0}
\definecolor{rev9color}{RGB}{128, 128, 128}
\definecolor{rev10color}{RGB}{102, 0, 204}

\def\PaperID{1862}
\def\confName{ICME}
\def\confYear{2026}

\usepackage{float}
\restylefloat{table}

\begin{document}

\title{ACME-YOLO: An Enhanced Small Object Detection Framework in UAVs}
\maketitle
\thispagestyle{empty}
\appendix

% ==============================
% 表1
% ==============================
\begin{table}[t]
  \caption{\textbf{Across-architecture study on VisDrone2019, UAVDT, TinyPerson, and DOTAv1.} Superscript \textsuperscript{\tiny +} denotes modified versions. Progressively deeper colors highlight improvements ($\Delta\geq0.005$), indicating larger gains over the unenhanced architecture.}
  \label{tab:across_arch}
  \centering
  % 适度调大列间距，避免内容挤压横线（原0.03pt过小，易遮挡）
  \setlength\tabcolsep{0.5pt}
  % 行高系数微调，既紧凑又不遮挡横线（原0.7→0.75）
  \renewcommand{\arraystretch}{0.75}
  % 字体行间距适配，减少垂直空白（原3pt→2.8pt）
  \fontsize{2.2pt}{2.8pt}\selectfont
  \resizebox{\columnwidth}{!}{
  \begin{tabular}{@{}>{\raggedright\arraybackslash}p{8em} 
                  @{}*{7}{>{\fontsize{2.2pt}{2.8pt}\selectfont}p{0.4cm}}
                  >{\fontsize{2.2pt}{2.8pt}\selectfont}p{0.4cm}
                  >{\fontsize{2.2pt}{2.8pt}\selectfont}p{0.4cm}
                  >{\fontsize{2.2pt}{2.8pt}\selectfont}p{0.4cm}
                  >{\fontsize{2.2pt}{2.8pt}\selectfont}p{0.18cm}
                  >{\fontsize{2.2pt}{2.8pt}\selectfont}p{0.4cm}@{}}
  \toprule[0.5pt]  % 明确横线粗细，避免过细被遮挡
  \textbf{Architecture}
  & \multicolumn{7}{c}{\fontsize{2.2pt}{2.8pt}\selectfont\textbf{Performance/mAP50}}
  & \multicolumn{5}{c}{\centering\fontsize{2.2pt}{2.8pt}\selectfont\textbf{Efficiency}} \\
  % 用\noalign缩小cmidrule上方间距，替代原[-2.2pt]（更稳定）
  \noalign{\vspace{-2pt}}
  \cmidrule[0.5pt](l{-0.3em}r{0.5em}){2-8}\cmidrule[0.5pt](l{-0.3em}r{0.5em}){10-13}
  & \multicolumn{2}{l}{\fontsize{2.2pt}{2.8pt}\selectfont\textbf{VisDrone2019}}
  & \multicolumn{2}{l}{\fontsize{2.2pt}{2.8pt}\selectfont\textbf{UAVDT}}
  & \multicolumn{2}{l}{\fontsize{2.2pt}{2.8pt}\selectfont\textbf{TinyPerson}}
  & \multicolumn{1}{l}{\fontsize{2.2pt}{2.8pt}\selectfont\textbf{DOTAv1}}
  & & & & & \\
  % 用\noalign缩小cmidrule上方间距，替代原[-2.2pt]
  \noalign{\vspace{-2pt}}
  \cmidrule[0.5pt](l{-0.3em}r{0.5em}){2-3}\cmidrule[0.5pt](l{-0.3em}r{0.5em}){4-5}\cmidrule[0.5pt](l{-0.3em}r{0.5em}){6-7}\cmidrule[0.5pt](l{-0.3em}r{0.5em}){8-8}
  & \cellcolor{blue!5}{\textbf{Val}}
  & \cellcolor{blue!5}{\textbf{Test}}
  & \cellcolor{blue!5}{\textbf{Val}}
  & \cellcolor{blue!5}{\textbf{Test}}
  & \cellcolor{blue!5}{\textbf{Val}}
  & \cellcolor{blue!5}{\textbf{Test}}
  & \cellcolor{blue!5}{\textbf{Val}}
  & \cellcolor{blue!5}{\textbf{Input}}
  & \cellcolor{blue!5}{\textbf{Params/M}}
  & \cellcolor{blue!5}{\textbf{GFLOPs}}
  & \cellcolor{blue!5}{\textbf{FPS}}
  & \cellcolor{blue!5}{\textbf{Size/MB}}\\
  % 核心：缩小midrule与数据的间距，且不遮挡横线
  \midrule[0.5pt]
  \noalign{\vspace{-2pt}}
  \textbf{YOLOv10-N}
  & \mapcellSI{0.328}{} & \mapcellSI{0.261}{} & \mapcellSI{0.845}{} & \mapcellSI{0.876}{} & \mapcellSI{0.225}{} & \mapcellSI{0.197}{} & \mapcellSI{0.422}{}
  & 640 & 2.7 & 8.4 & 187 & 5.9 \\
  \textbf{ACME-YOLO-N}
  & \mapcellSI{0.347}{0.019} & \mapcellSI{0.285}{0.024} & \mapcellSI{0.858}{0.013} & \mapcellSI{0.890}{0.014} & \mapcellSI{0.243}{0.018} & \mapcellSI{0.212}{0.015} & \mapcellSI{0.441}{0.019}
  & 640 & \effcell{2.7}{3.0}{+11.1} & \effcell{8.4}{8.7}{+3.6} & \fpscell{182} & \effcell{5.9}{6.3}{+6.8} \\
  % 核心：缩小bottomrule与数据的间距，且不遮挡横线
  \noalign{\vspace{-2pt}}
  \bottomrule[0.5pt]
  \end{tabular}}
  \vspace{-0.2cm}
\end{table}

% ==============================
% 表2
% ==============================


% ==============================
% 回复正文（所有 Q/A 完全左对齐，无缩进）
% ==============================

\noindent\textcolor{blue}{\textbf{$\triangleright$ ALL Reviewers:}}

\noindent We appreciate all reviewers for their valuable feedback and constructive suggestions. We have carefully addressed all concerns and revised the manuscript accordingly.

\noindent\textbf{Q1:} Concern on Incremental Novelty / Engineering Combination\\
\noindent\textbf{A:} We agree that attention, multi-scale dilation, and upsampling have been explored before. Our contribution is not on novel basic operators, but a \textbf{constraint-aware structural co-design} tailored for lightweight UAV small-object detection.

\noindent CCRE integrates iAFF into YOLOv10’s C2f backbone with a dual-path residual structure and learnable scaling factor $\alpha$, enhancing shallow P2/P3 discriminability under tight parameter constraints. To our knowledge, such a lightweight residual-attention design for UAV small-object detection has not been studied.

\noindent ADRF-Fusion uses parallel depthwise dilated branches with Softmax-normalized weights for adaptive multi-scale fusion with negligible extra complexity, avoiding heavy fusion modules.

\noindent Notably, these modules form a unified \textbf{enhancement–fusion–restoration pipeline} that addresses mutually coupled bottlenecks, rather than being independently plugged in. Ablation results verify their complementary effects.

\noindent Thus, our work represents a \textbf{structured, lightweight, synergistic redesign} under practical deployment constraints.

\vspace{0.5em}

% ----------------------------
% Reviewer 6
% ----------------------------
\noindent\textcolor{rev6color}{\textbf{$\triangleright$ Reviewer\#6.}}

\noindent\textbf{Q1:} Concern on Limited Performance Gain and Robustness\\
\noindent\textbf{A:} Although the absolute gain (+1.9 mAP50 on VisDrone2019) is moderate, it is achieved with only ~0.3M parameter increase and negligible FLOPs growth, representing a favorable accuracy–efficiency trade-off for deployment scenarios. Improvements are consistent across VisDrone and UAVDT (+1.53 mAP50), suggesting cross-dataset generalization rather than dataset-specific tuning.

\noindent\textbf{Q2:} Concern on Real-Time Claim Without FPS / Edge Hardware Evidence\\
\noindent\textbf{A:} We acknowledge that Params/FLOPs alone are insufficient to support real-time claims. We will provide explicit end-to-end FPS comparisons to ensure all deployment claims are strictly empirically supported.

\vspace{0.5em}

% ----------------------------
% Reviewer 7
% ----------------------------
\noindent\textcolor{rev7color}{\textbf{$\triangleright$ Reviewer\#7.}}
\noindent (All comments from Reviewer 7 have been addressed in the earlier responses.)

\vspace{0.5em}
\begin{table}[t]
  \caption{\textbf{TIDE error decomposition.} Values are reported for YOLOv10-N and ACME-YOLO-N.}
  \label{tab:tide_nano}
  \centering
  % 1. 适度增大列间距，避免内容挤压横线
  \setlength\tabcolsep{2pt}
  % 2. 行高系数调至适中，既紧凑又不遮挡
  \renewcommand{\arraystretch}{1.1}
  % 3. 字体行间距与字号匹配，避免垂直挤压
  \fontsize{4pt}{4.8pt}\selectfont
  % 移除resizebox的强制拉伸，改用固定宽度更可控（核心解决遮挡）
  \begin{tabular}{@{}>{\raggedright\arraybackslash}p{1.2cm} 
                  *{4}{>{\raggedright\arraybackslash}p{0.5cm}}  % 前4列宽度不变
                  >{\raggedright\arraybackslash}p{0.8cm}       % FalseNeg列调大（0.5→0.8，更宽松）
                  >{\raggedright\arraybackslash}p{0.95cm}@{}} % APs列宽度不变
    \toprule[0.6pt]  % 稍加粗横线，避免被背景色/内容遮挡
    \textbf{Architecture}
    & \cellcolor{blue!5}{\fontsize{4.5pt}{5pt}\selectfont\textbf{Loc}}
    & \cellcolor{blue!5}{\fontsize{4.5pt}{5pt}\selectfont\textbf{Miss}}
    & \cellcolor{blue!5}{\fontsize{4.5pt}{5pt}\selectfont\textbf{Both}}
    & \cellcolor{blue!5}{\fontsize{4.5pt}{5pt}\selectfont\textbf{FalsePos}}
    & \cellcolor{blue!5}{\fontsize{4.5pt}{5pt}\selectfont\textbf{FalseNeg}}
    & \cellcolor{blue!5}{\fontsize{4.5pt}{5pt}\selectfont\textbf{APs}}\\
    % 4. 用\noalign缩小横线与数据的间距（最安全的方式）
    \midrule[0.6pt]
    \noalign{\vspace{-1.5pt}}
    \textbf{YOLOv10-N}
    & 3.54 & 6.44 & 0.25 & 13.21 & 26.46 & 7.69 \\
    \textbf{ACME-YOLO-N}
    & 3.30\tidedown & 6.15\tidedown & 0.23\tidedown & 13.08\tidedown & 24.37\tidedown & 8.84\tideupAB{1.15}{} \\
    \noalign{\vspace{-1.5pt}}
    \bottomrule[0.6pt]
  \end{tabular}
  \vspace{-0.2cm}
\end{table}
% ----------------------------
% Reviewer 8
% ----------------------------
\noindent\textcolor{rev8color}{\textbf{$\triangleright$ Reviewer\#8.}}

\noindent\textbf{Q1:} Concern on the rationality of "Adaptive" in ADRF-Fusion\\
\noindent\textbf{A:} The “Adaptive” refers to data-driven, task-adaptive weights learned end-to-end via Softmax, rather than manually fixed fusion weights. Although global, these weights automatically adapt to multi-scale feature importance for small object detection.

\noindent\textbf{Q2:} Concern on missing FPS/latency and deployment results\\
\noindent\textbf{A:} As shown in \cref{tab:across_arch}, we have reported FPS at 640×640, batch-size 1 on an A100 GPU. We have not performed official ONNX/TensorRT deployment yet, but will add standard deployment evaluation in the revised version.

\noindent\textbf{Q3:} Concern on insufficient small-object evaluation metrics\\
\noindent\textbf{A:} As presented in \cref{tab:tide_nano}, we have provided APs (small-object AP) to evaluate small object detection performance.

\noindent\textbf{Q4:} Concern on inconsistent training settings among YOLO variants\\
\noindent\textbf{A:} Thank you for the careful comment. All compared YOLO variants are trained with the same training strategy, input size, epochs, and augmentations. We will add a clear clarification in the revised manuscript.

\vspace{0.5em}

% ----------------------------
% Reviewer 9
% ----------------------------
\noindent\textcolor{rev9color}{\textbf{$\triangleright$ Reviewer\#9.}}

\vspace{0.5em}

% ----------------------------
% Reviewer 10
% ----------------------------
\noindent\textcolor{rev10color}{\textbf{$\triangleright$ Reviewer\#10.}}

\noindent\textbf{Q1:} Concern on insufficient efficiency evidence with increased Params/GFLOPs\\
\noindent\textbf{A:} As shown in \cref{tab:across_arch}, we have provided FPS measurements on an A100 GPU, which together with params and GFLOPs fully demonstrate the accuracy-efficiency trade-off of our method.\\
\noindent\textbf{Q2:} Typos, small figure fonts, and formatting issues\\
\noindent\textbf{A:} Thanks for the careful check. We will fix all typos, adjust fonts, and polish formatting in the revised manuscript.
\end{document}
